%% ============================================================
%% webservice.tex — 2025년 겨울·준비단계 구현 내용 (배포된 웹 서비스, 라이브 캡처)
%% 섹션 1(팀 구성) 바로 뒤에 배치 — 현장 개발(1~11일차)의 기반이 된 준비단계 산출물
%% 출처: https://dknife.github.io/ARES2 (2026-06-27 캡처)
%% ============================================================
\section{2025년 겨울·준비단계 구현 내용 --- 배포된 웹 서비스}
6월 현장 착수 이전, \textbf{2025년 겨울부터 준비단계}에 구현해 GitHub Pages로 배포한
ARES 웹 서비스(\href{https://dknife.github.io/ARES2}{\code{dknife.github.io/ARES2}})의
주요 기능을 정리한다. \textbf{이후 1$\sim$11일차의 현장 개발은 이 준비단계 결과물을
기반으로} 진행됐다. 설치 없이 브라우저만으로 학습 콘텐츠 열람, 블록 코딩, 3D
시뮬레이션, AI 도움, 기기 모니터링이 가능하며, 블록 1개가 명령 문자열로 변환되어
BLE로 로버에 전달되는 전체 경로(\ref{app:code})를 실물 기기 없이도 시뮬레이션으로
체험할 수 있다.

\subsection{랜딩 페이지와 학습 관제실}
랜딩 페이지는 제품 정체성(``블록 코딩으로 화성을 탐사하다'')을 제시하고, ``서비스
시작''으로 학습 관제실에 진입한다. 관제실은 \textbf{12차시 $\times$ 4미션 = 48개 미션}의
커리큘럼(‘화성 탐사선 코딩 여정’)을 차시별 카드와 ‘차시 한눈에 보기’ 표로 안내한다.

\begin{figure}[H]\centering
  \begin{subfigure}{0.49\linewidth}\centering
    \figbox{webcaps/web_landing.png}
    \caption{랜딩 페이지 — 제품 소개 히어로}
  \end{subfigure}\hfill
  \begin{subfigure}{0.49\linewidth}\centering
    \figbox{webcaps/web_overview.png}
    \caption{학습 관제실 — 12차시 커리큘럼$\cdot$차시 한눈에 보기}
  \end{subfigure}
  \caption{배포 웹 서비스 — 랜딩 페이지와 학습 관제실(차시$\cdot$미션 선택)}
\end{figure}

\subsection{블록 코딩과 3D 시뮬레이션}
미션을 고르면 Google Blockly 기반 블록 코딩 화면이 열린다. 좌측 툴박스는
서보$\cdot$DC 모터$\cdot$LED$\cdot$디스플레이$\cdot$소리$\cdot$발사$\cdot$센서$\cdot$시간$\cdot$제어$\cdot$변수$\cdot$수학$\cdot$함수
등 도메인별 카테고리로 구성된다. ‘예제’를 불러오면(아래는 \emph{장애물 피하기})
거리 센서 분기$\cdot$DC 모터 주행$\cdot$LED$\cdot$부저 블록이 조합된 완성 프로그램이 채워진다.
‘시뮬레이션’을 열면 탐사로봇 ‘알비’가 3D로 등장해, 작성한 블록이 기기 없이도
어떻게 동작하는지 보여 준다(WebGL~\cite{webgl, threejs}).

\begin{figure}[H]\centering
  \begin{subfigure}{0.49\linewidth}\centering
    \figbox{webcaps/web_blockcoding.png}
    \caption{블록 코딩 — ‘장애물 피하기’ 예제(센서$\cdot$DC$\cdot$LED$\cdot$부저)}
  \end{subfigure}\hfill
  \begin{subfigure}{0.49\linewidth}\centering
    \figbox{webcaps/web_simulation.png}
    \caption{3D 시뮬레이션 — 탐사로봇 ‘알비’ 동작 미리보기}
  \end{subfigure}
  \caption{배포 웹 서비스 — 블록 코딩 워크스페이스와 WebGL 3D 시뮬레이션}
\end{figure}

\subsection{시뮬레이션 모델 — 4종 전환}
시뮬레이션 화면 좌측 상단의 \textbf{모델 선택 드롭다운}으로 \textbf{4종의 시뮬레이션
모델}을 전환할 수 있다. 같은 블록 코드라도 선택한 모델에 따라 서로 다른 3D 장면과
동작이 연출되어, 미션 맥락(로봇$\cdot$신호등$\cdot$발사$\cdot$주행)에 맞는 학습이 가능하다.
네 모델을 차례로 선택해 각각 다른 시뮬레이션이 이뤄짐을 확인했다.

\begin{itemize}
  \item \textbf{알비와 함께} — 탐사로봇 ‘알비’와 LED(눈$\cdot$가슴) 제어 중심의 입문 장면.
  \item \textbf{우주 신호등} — 신호등 슬롯$\cdot$램프와 가위바위보 손동작(랜덤 함수 학습).
  \item \textbf{발사대} — 로켓 발사 시퀀스(레이더$\cdot$로켓$\cdot$LED), 점화$\cdot$발사 연출.
  \item \textbf{로버} — DC 모터 주행과 거리 센서$\cdot$장애물 회피(격자 바닥의 장애물 박스).
\end{itemize}

\begin{figure}[H]\centering
  \begin{subfigure}{0.49\linewidth}\centering
    \figbox{webcaps/sim_albi.png}
    \caption{알비와 함께 — 로봇$\cdot$LED}
  \end{subfigure}\hfill
  \begin{subfigure}{0.49\linewidth}\centering
    \figbox{webcaps/sim_traffic.png}
    \caption{우주 신호등 — 신호등$\cdot$가위바위보}
  \end{subfigure}

  \medskip
  \begin{subfigure}{0.49\linewidth}\centering
    \figbox{webcaps/sim_launchpad.png}
    \caption{발사대 — 로켓 발사 시퀀스}
  \end{subfigure}\hfill
  \begin{subfigure}{0.49\linewidth}\centering
    \figbox{webcaps/sim_rover.png}
    \caption{로버 — DC 주행$\cdot$장애물 회피}
  \end{subfigure}
  \caption{시뮬레이션 모델 4종 — 드롭다운으로 전환하면 서로 다른 3D 장면$\cdot$동작이 연출된다}
\end{figure}

\subsection{AI 도움과 기기 모니터링}
‘AI 도움’ 패널은 하고 싶은 일을 자연어로 적으면 블록 구성을 제안한다. ‘점검’
화면(ARES 미션 컨트롤 대시보드)은 BLE 연결 상태, 센서 텔레메트리(거리$\cdot$자기장$\cdot$온도),
시스템 자원(메모리$\cdot$가동 시간)과 수동 제어(전진$\cdot$정지$\cdot$LED$\cdot$회전$\cdot$소리)를
한눈에 보여 준다.

\begin{figure}[H]\centering
  \begin{subfigure}{0.49\linewidth}\centering
    \figbox{webcaps/web_ai.png}
    \caption{AI 도움 — 자연어로 블록 구성 제안}
  \end{subfigure}\hfill
  \begin{subfigure}{0.49\linewidth}\centering
    \figbox{webcaps/web_dashboard.png}
    \caption{점검 대시보드 — 연결$\cdot$센서$\cdot$시스템$\cdot$수동 제어}
  \end{subfigure}
  \caption{배포 웹 서비스 — AI 도움 패널과 기기 모니터링 대시보드}
\end{figure}

\begin{notebox}[캡처 환경]
위 화면은 \code{dknife.github.io/ARES2}를 데스크톱 브라우저로 구동해 2026-06-27에
캡처했다. 블루투스(BLE)는 사용자 제스처$\cdot$실기기가 필요해 연결 상태는 ‘연결 끊김’으로
표시되며, 블록 코딩$\cdot$3D 시뮬레이션$\cdot$학습 콘텐츠는 기기 없이도 동작한다.
\end{notebox}
